
%% ===================================================================== C
\section{Part C --- Textbook building blocks}

\begin{qbox}{C1.\ $M_{X}M_{1} = M_{X}$}{\likely\ \etmtag{2.15}}
Show $M_{X}M_{1} = M_{1}M_{X} = M_{X}$ when $X = [X_{1}\ X_{2}]$.
\end{qbox}
\begin{solbox}
$X_{1}\subset X\Rightarrow P_{X}X_{1}=X_{1}\Rightarrow P_{X}P_{1}=P_{X}X_{1}(X_{1}\xT X_{1})^{-1}X_{1}\xT=P_{1}$. Then
\[
  M_{X}M_{1}=(I-P_{X})(I-P_{1})=I-P_{1}-P_{X}+P_{X}P_{1}=I-P_{X}=M_{X}.
\]
Transposing gives $M_{1}M_{X}=M_{X}$.
\begin{why}\Why removing the $X_{1}$-part first and then everything in $X$ is the same as removing everything in $X$,
because $S(X_{1})\subset S(X)$.\end{why}
\end{solbox}

\begin{qbox}{C2.\ $P_{X} - P_{1} = P_{M_{1}X_{2}}$}{\must\ \etmtag{2.19}}
Show $P_{X} - P_{1} = P_{M_{1}X_{2}}$. Use it to show
$y\xT(P_{X} - P_{1})y = y\xT M_{1}X_{2}(X_{2}\xT M_{1}X_{2})^{-1}X_{2}\xT M_{1}y$.
\end{qbox}
\begin{solbox}
Let $Z=M_{1}X_{2}$.
\step{1. $S(X)=S(X_{1})+S(Z)$.} $X_{1}c_{1}+X_{2}c_{2}=X_{1}\big(c_{1}+(X_{1}\xT X_{1})^{-1}X_{1}\xT X_{2}c_{2}\big)+Zc_{2}$,
and conversely any element of $S(X_{1})+S(Z)$ is in $S(X)$.
\step{2. Orthogonal.} $X_{1}\xT Z=(M_{1}X_{1})\xT X_{2}=O$, so $P_{1}Z=O$, hence $P_{1}P_{Z}=O=P_{Z}P_{1}$.
\step{3. Sum of projections.} $Q=P_{1}+P_{Z}$ is symmetric and $Q^{2}=P_{1}+P_{Z}=Q$.
On $v=X_{1}d+Zc\in S(X)$: $Qv=v$. On $w\perp S(X)$: $X_{1}\xT w=0$ gives $P_{1}w=0$, and
$Z\xT w=X_{2}\xT w-X_{2}\xT P_{1}w=0$ gives $P_{Z}w=0$. So $Q=P_{X}$.
\step{4. Quadratic form.} $Z\xT Z=X_{2}\xT M_{1}X_{2}$, $Z\xT y=X_{2}\xT M_{1}y$, so
$y\xT(P_{X}-P_{1})y=y\xT Z(Z\xT Z)^{-1}Z\xT y=y\xT M_{1}X_{2}(X_{2}\xT M_{1}X_{2})^{-1}X_{2}\xT M_{1}y$.
\begin{why}\Why what $X$ explains beyond $X_{1}$ is what the \emph{new} part of $X_{2}$, $M_{1}X_{2}$, explains. Also
$y\xT M_{1}y-y\xT M_{X}y=y\xT(P_{X}-P_{1})y$, which is $RSSR-USSR$ for $H_{0}:\beta_{2}=0$: the $F$ numerator.\end{why}
\end{solbox}

\begin{qbox}{C3.\ Expected residual sum of squares}{\likely\ \etmtag{3.14}}
Prove $\E(\hat u\xT\hat u) = (n-k)\sigma^{2}$ using the trace.
\end{qbox}
\begin{solbox}
$\hat u=M_{X}y=M_{X}(X\beta+u)=M_{X}u$, so $\hat u\xT\hat u=u\xT M_{X}u$. A scalar equals its trace and
$\mathrm{tr}(AB)=\mathrm{tr}(BA)$:
\[
  \E(u\xT M_{X}u)=\E\,\mathrm{tr}(M_{X}uu\xT)=\mathrm{tr}(M_{X}\,\sigma^{2}I)=\sigma^{2}\,\mathrm{tr}(M_{X}),
\]
and $\mathrm{tr}(P_{X})=\mathrm{tr}\big((X\xT X)^{-1}X\xT X\big)=k$, so $\mathrm{tr}(M_{X})=n-k$.
\begin{why}\Why the $n$-dimensional error loses $k$ dimensions to the fit; $n-k$ dimensions of noise remain, so divide by
$n-k$ for an unbiased $s^{2}$.\end{why}
\end{solbox}

\begin{qbox}{C4.\ Orthogonal regressors}{\opt\ \etmtag{3.15}}
For $y = X\beta + u$ and $y = X\beta + Z\gamma + u$ with $X\xT Z = O$, show the OLS estimates of $\beta$ coincide.
\end{qbox}
\begin{solbox}
FWL: the long-regression estimate is $(X\xT M_{Z}X)^{-1}X\xT M_{Z}y$. If $X\xT Z=O$ then $M_{Z}X=X-Z(Z\xT Z)^{-1}Z\xT X=X$, so it
equals $(X\xT X)^{-1}X\xT y$.
\begin{why}\Why cleaning $X$ of $Z$ changes nothing when they share nothing; omitting an orthogonal variable causes no bias.\end{why}
\end{solbox}

\begin{qbox}{C5.\ Restricted vs unrestricted estimates}{\likely\ \etmtag{4.10}}
Show that the regression of $y$ on $X_{1}$ alone gives
$\tilde\beta_{1}=\hat\beta_{1}+(X_{1}\xT X_{1})^{-1}X_{1}\xT X_{2}\hat\beta_{2}$, where $\hat\beta_{1},\hat\beta_{2}$ are from the long regression.
\end{qbox}
\begin{solbox}
The first block of the long regression's normal equations is $X_{1}\xT X_{1}\hat\beta_{1}+X_{1}\xT X_{2}\hat\beta_{2}=X_{1}\xT y$. Multiply by
$(X_{1}\xT X_{1})^{-1}$:
$\hat\beta_{1}+(X_{1}\xT X_{1})^{-1}X_{1}\xT X_{2}\hat\beta_{2}=(X_{1}\xT X_{1})^{-1}X_{1}\xT y=\tilde\beta_{1}$.
\begin{why}\Why when $X_{2}$ is dropped its effect $X_{2}\hat\beta_{2}$ is handed to $X_{1}$ through the auxiliary coefficients
$(X_{1}\xT X_{1})^{-1}X_{1}\xT X_{2}$. Taking expectations gives the omitted-variable bias.\end{why}
\end{solbox}

\begin{qbox}{C6.\ Chow test when a subsample is short}{\must\ \etmtag{4.12}\ \examtag{=A1(c)}}
When one subsample has $n_{2}\le k$ observations, show the Chow statistic becomes
$F=\dfrac{(RSSR-SSR_{1})/n_{2}}{SSR_{1}/(n_{1}-k)}$ and that numerator and denominator have these degrees of freedom.
\end{qbox}
\begin{solbox}
\step{Dummies.} Unrestricted model: $y=X\beta+D\delta+u$ with $D=\begin{bmatrix}O\\I_{n_{2}}\end{bmatrix}$ (one dummy per
subsample-2 observation). Each dummy fits its observation exactly, so the RSS equals $SSR_{1}$.
\step{Count.} Restricted model: $\delta=0$ (one common line), RSS $=RSSR$, $q=n_{2}$. Unrestricted parameters $k+n_{2}$, so the
denominator df is $n-k-n_{2}=n_{1}-k$:
$F=\dfrac{(RSSR-SSR_{1})/n_{2}}{SSR_{1}/(n_{1}-k)}\sim F(n_{2},n_{1}-k)$.
\begin{why}\Why with $n_{2}\le k$ the second subsample cannot estimate its own $\beta$; only $n_{2}$ pieces of information are
available, so only $n_{2}$ restrictions are tested (the ``predictive failure'' test).\end{why}
\end{solbox}

%% ===================================================================== D
\section{Part D --- New questions on Lessons 1--8}

\begin{qbox}{D1.\ Numerical: $t$, overall $F$, joint $F$}{\likely\ \newtag}
A regression of $y$ on a constant and $x_{2},x_{3},x_{4}$ with $n=50$ gives $b_{2}=0.5\ (0.2)$, $b_{3}=-1.2\ (0.8)$,
$b_{4}=0.9\ (0.3)$ (standard errors in parentheses), $ESS=80$, $RSS=40$.
\begin{flatlist}
  \item[(a)] Test each slope at 5\%. Compute $R^{2}$ and $\bar R^{2}$.
  \item[(b)] Test $H_{0}:\beta_{2}=\beta_{3}=\beta_{4}=0$.
  \item[(c)] Two more regressors are added and $RSS$ falls to $36$. Test their joint significance.
  \item[(d)] How would you test $\beta_{2}+\beta_{4}=1$ with (i) a $t$ test and (ii) a restricted regression?
        What is missing from the output for (i)?
\end{flatlist}
\end{qbox}
\begin{solbox}
\step{(a)} $n-k=46$, $t_{0.025}(46)\approx2.01$: $t_{2}=2.5$ (significant), $t_{3}=-1.5$ (not), $t_{4}=3.0$ (significant).
$TSS=120$, $R^{2}=80/120=0.667$, $\bar R^{2}=1-\frac{40/46}{120/49}=0.645$.
\step{(b)} $F=\frac{ESS/3}{RSS/46}=\frac{26.67}{0.870}=30.7>F_{0.05}(3,46)\approx2.81$: reject.
\step{(c)} $F=\frac{(40-36)/2}{36/44}=2.44<F_{0.05}(2,44)\approx3.21$: do not reject.
\step{(d)} (i) $t=\frac{b_{2}+b_{4}-1}{\sqrt{\Var(b_{2})+\Var(b_{4})+2\Cov(b_{2},b_{4})}}=\frac{0.4}{\sqrt{0.13+2\Cov(b_{2},b_{4})}}$;
the covariance $\Cov(b_{2},b_{4})$ is not reported. (ii) Substitute $\beta_{4}=1-\beta_{2}$: regress $y-x_{4}$ on a constant, $x_{2}-x_{4}$ and $x_{3}$;
then $F=\frac{(RSSR-40)/1}{40/46}\sim F(1,46)$, equal to $t^{2}$ from (i).
\begin{why}\Why a linear combination has variance $w\xT\Var(b)w$, which needs the off-diagonal term; reparametrizing turns the restriction into an
exclusion so the RSS comparison does the work.\end{why}
\end{solbox}

\begin{qbox}{D2.\ Chow test with enough observations}{\likely\ \newtag}
A model with $k=3$ is estimated on $n=50$ observations, $RSS=60$. On the first $n_{1}=30$, $SSR_{1}=20$; on the last $n_{2}=20$,
$SSR_{2}=25$. Test parameter stability. State the assumption on the error variances. How does the test change if $n_{2}=2$?
\end{qbox}
\begin{solbox}
$USSR=20+25=45$, $n-2k=44$, $q=k=3$: $F=\frac{(60-45)/3}{45/44}=4.89>F_{0.05}(3,44)\approx2.82$.
\Ans{Reject stability. The test assumes the same $\sigma^{2}$ in both subsamples (and independent normal errors). If $n_{2}=2<k$:
$F=\frac{(RSSR-SSR_{1})/2}{SSR_{1}/(n_{1}-3)}$ with $n_{1}=48$, i.e.\ $F(2,45)$ (C6).}
\end{solbox}

\begin{qbox}{D3.\ Gauss--Markov}{\likely\ \newtag}
In $y=X\beta+u$, $\E(u)=0$, $\Var(u)=\sigma^{2}I$, $X$ fixed of full rank, let $b^{*}=Cy$ be any linear unbiased estimator with
$C=(X\xT X)^{-1}X\xT+A$. Show unbiasedness requires $AX=O$ and $\Var(b^{*})=\Var(b)+\sigma^{2}AA\xT$. Conclude $w\xT b$ is BLUE for $w\xT\beta$.
\end{qbox}
\begin{solbox}
\step{Unbiased.} $\E(b^{*})=CX\beta=\beta+AX\beta$ for all $\beta$ $\iff$ $AX=O$.
\step{Variance.} $b^{*}-\beta=Cu$, so $\Var(b^{*})=\sigma^{2}CC\xT$ and
$CC\xT=(X\xT X)^{-1}+\underbrace{(X\xT X)^{-1}(AX)\xT+AX(X\xT X)^{-1}}_{=O}+AA\xT$, giving $\Var(b^{*})=\Var(b)+\sigma^{2}AA\xT$.
\step{BLUE.} $\Var(w\xT b^{*})-\Var(w\xT b)=\sigma^{2}\|A\xT w\|^{2}\ge0$.
\begin{why}\Why any other linear unbiased estimator is OLS plus a piece $Au$ uncorrelated with OLS (cross terms vanish since $AX=O$);
adding uncorrelated noise can only add variance.\end{why}
\end{solbox}

\begin{qbox}{D4.\ Quadratic forms, $s^{2}$ and independence}{\likely\ \newtag}
\begin{flatlist}
  \item[(a)] Using the spectral decomposition, show that $z\sim N(0,I_{n})$ and $A$ symmetric idempotent of rank $r$ imply $z\xT Az\sim\chi^{2}(r)$.
  \item[(b)] Show $(n-k)s^{2}/\sigma^{2}=u\xT M_{X}u/\sigma^{2}\sim\chi^{2}(n-k)$.
  \item[(c)] Show $b$ and $\hat u$ are independent, hence $b$ and $s^{2}$ are independent.
  \item[(d)] Conclude $(b_{j}-\beta_{j})/SE(b_{j})\sim t(n-k)$.
\end{flatlist}
\end{qbox}
\begin{solbox}
\step{(a)} $A=Q\Lambda Q\xT$, $Q$ orthogonal. $\lambda^{2}=\lambda$ so eigenvalues are $0$ or $1$, with $r=\mathrm{tr}A$ ones. $v=Q\xT z\sim N(0,I)$ and
$z\xT Az=v\xT\Lambda v=\sum_{i=1}^{r}v_{i}^{2}\sim\chi^{2}(r)$.
\step{(b)} $(n-k)s^{2}=u\xT M_{X}u$ (C3), $u/\sigma\sim N(0,I)$, $M_{X}$ idempotent of rank $n-k$, so (a) applies.
\step{(c)} $b-\beta=(X\xT X)^{-1}X\xT u$ and $\hat u=M_{X}u$ are jointly normal, and
$\Cov(b,\hat u)=\sigma^{2}(X\xT X)^{-1}X\xT M_{X}=O$. Jointly normal and uncorrelated $\Rightarrow$ independent; $s^{2}$ depends only on $\hat u$.
\step{(d)} With $c_{jj}=[(X\xT X)^{-1}]_{jj}$:
$\frac{b_{j}-\beta_{j}}{s\sqrt{c_{jj}}}=\frac{(b_{j}-\beta_{j})/(\sigma\sqrt{c_{jj}})}{\sqrt{[(n-k)s^{2}/\sigma^{2}]/(n-k)}}=\frac{N(0,1)}{\sqrt{\chi^{2}(n-k)/(n-k)}}\sim t(n-k)$.
\begin{why}\Why $b$ uses the part of $u$ in $S(X)$ and $\hat u$ the part in $S(X)^{\perp}$; orthogonal pieces of a spherical normal vector are independent.\end{why}
\end{solbox}

\begin{qbox}{D5.\ The $F$ test and restricted least squares}{\must\ \newtag}
For $H_{0}:R\beta=r$ ($q$ restrictions):
\begin{flatlist}
  \item[(a)] Show $(Rb-r)\xT[\sigma^{2}R(X\xT X)^{-1}R\xT]^{-1}(Rb-r)\sim\chi^{2}(q)$ under $H_{0}$.
  \item[(b)] Minimize $(y-X\beta)\xT(y-X\beta)$ s.t.\ $R\beta=r$. Show $b_{*}=b-(X\xT X)^{-1}R\xT[R(X\xT X)^{-1}R\xT]^{-1}(Rb-r)$ and find $\lambda_{*}$.
  \item[(c)] Show $e_{*}\xT e_{*}-e\xT e=(Rb-r)\xT[R(X\xT X)^{-1}R\xT]^{-1}(Rb-r)$, so $F=\frac{(RSSR-USSR)/q}{USSR/(n-k)}\sim F(q,n-k)$.
  \item[(d)] Find $\E(b_{*})$, $\Var(b_{*})$ when $H_{0}$ is true and false. Show $\Var(b)-\Var(b_{*})$ is p.s.d.
  \item[(e)] With $\sigma^{2}$ known, show the Wald statistic (from $Rb-r$) and LM statistic (from $\lambda_{*}$) are identical.
  \item[(f)] For $R=[O\ \ I_{k_{2}}]$, $r=0$, show the numerator is $y\xT M_{1}X_{2}(X_{2}\xT M_{1}X_{2})^{-1}X_{2}\xT M_{1}y$.
\end{flatlist}
\end{qbox}
\begin{solbox}
Let $W\equiv R(X\xT X)^{-1}R\xT$ ($q\times q$, invertible).
\step{(a)} Under $H_{0}$, $Rb-r=R(X\xT X)^{-1}X\xT u\sim N(0,\sigma^{2}W)$. For $w\sim N(0,V)$, $V^{-1/2}w\sim N(0,I_{q})$, so $w\xT V^{-1}w\sim\chi^{2}(q)$.
\step{(b)} $\mathcal L=(y-X\beta)\xT(y-X\beta)+2\lambda\xT(R\beta-r)$. FOC: $-2X\xT y+2X\xT X\beta+2R\xT\lambda=0\Rightarrow b_{*}=b-(X\xT X)^{-1}R\xT\lambda$.
Imposing $Rb_{*}=r$: $Rb-W\lambda=r$, so
$\lambda_{*}=W^{-1}(Rb-r)$ and $b_{*}=b-(X\xT X)^{-1}R\xT W^{-1}(Rb-r)$.
\step{(c)} $e_{*}=e+X(b-b_{*})$ and $X\xT e=0$, so
$e_{*}\xT e_{*}=e\xT e+(b-b_{*})\xT X\xT X(b-b_{*})=e\xT e+(Rb-r)\xT W^{-1}(Rb-r)$.
Divided by $\sigma^{2}$, the gap is $\chi^{2}(q)$ by (a) and depends on $b$ only; $USSR/\sigma^{2}\sim\chi^{2}(n-k)$ depends on $e$ only; independent (D4c). So
$F=\frac{(RSSR-USSR)/q}{USSR/(n-k)}\sim F(q,n-k)$ and $\sigma^{2}$ cancels.
\begin{why}\Why imposing a restriction can only raise the RSS; the gap is how far $Rb$ is from $r$, in units of its own variance.\end{why}
\step{(d)} With $G=(X\xT X)^{-1}R\xT W^{-1}R(X\xT X)^{-1}$:
$\E(b_{*})=\beta-(X\xT X)^{-1}R\xT W^{-1}(R\beta-r)$ and $\Var(b_{*})=\sigma^{2}[(X\xT X)^{-1}-G]$ (same whether or not $H_{0}$ holds, since $r$ is fixed).
Under $H_{0}$ the bias vanishes; otherwise $b_{*}$ is biased. $\Var(b)-\Var(b_{*})=\sigma^{2}G=\sigma^{2}H\xT W^{-1}H$ with $H=R(X\xT X)^{-1}$: p.s.d.
\begin{why}\Why a correct restriction is free information (lower variance, no bias); a false one still lowers variance but adds bias.\end{why}
\step{(e)} Wald: $(Rb-r)\xT(\sigma^{2}W)^{-1}(Rb-r)$. LM: $\Var(\lambda_{*})=W^{-1}(\sigma^{2}W)W^{-1}=\sigma^{2}W^{-1}$, so
$LM=\lambda_{*}\xT(\sigma^{2}W^{-1})^{-1}\lambda_{*}=(Rb-r)\xT W^{-1}WW^{-1}(Rb-r)/\sigma^{2}$ = Wald.
\step{(f)} $Rb=b_{2}$ and $W=(X_{2}\xT M_{1}X_{2})^{-1}$ (B1), so the numerator is $b_{2}\xT(X_{2}\xT M_{1}X_{2})b_{2}$. With the FWL $b_{2}$ this is
$y\xT M_{1}X_{2}(X_{2}\xT M_{1}X_{2})^{-1}X_{2}\xT M_{1}y=y\xT(P_{X}-P_{1})y=RSSR-USSR$ (C2).
\end{solbox}

\begin{qbox}{D6.\ Omitted and irrelevant variables}{\likely\ \newtag}
True model $y=X_{1}\beta_{1}+X_{2}\beta_{2}+u$, $u\sim N(0,\sigma^{2}I)$.
\begin{flatlist}
  \item[(a)] Estimate $y=X_{1}\beta_{1}+u$. Derive the distribution of $\tilde\beta_{1}$. When is it unbiased?
  \item[(b)] Show $\Var(\hat\beta_{1})-\Var(\tilde\beta_{1})$ is p.s.d.
  \item[(c)] Compare $MSE(\tilde\beta_{1})$ with $MSE(\hat\beta_{1})$. When does omitting $X_{2}$ win?
  \item[(d)] Now the truth is $y=X_{1}\beta_{1}+u$ but you include $X_{2}$. Show the estimator is unbiased with larger variance, and $s^{2}$ is unbiased.
\end{flatlist}
\end{qbox}
\begin{solbox}
Let $a=(X_{1}\xT X_{1})^{-1}X_{1}\xT X_{2}$.
\step{(a)} $\tilde\beta_{1}=\beta_{1}+a\beta_{2}+(X_{1}\xT X_{1})^{-1}X_{1}\xT u\sim N\big(\beta_{1}+a\beta_{2},\sigma^{2}(X_{1}\xT X_{1})^{-1}\big)$;
unbiased iff $a\beta_{2}=0$ (e.g.\ $\beta_{2}=0$ or $X_{1}\xT X_{2}=O$).
\step{(b)} $\Var(\hat\beta_{1})=\sigma^{2}(X_{1}\xT M_{2}X_{1})^{-1}$ and $X_{1}\xT X_{1}-X_{1}\xT M_{2}X_{1}=X_{1}\xT P_{2}X_{1}\ge0$. For p.d.\ matrices $A\ge B\Rightarrow B^{-1}\ge A^{-1}$, so
$(X_{1}\xT M_{2}X_{1})^{-1}\ge(X_{1}\xT X_{1})^{-1}$.
\step{(c)} By C5 $\tilde\beta_{1}=\hat\beta_{1}+a\hat\beta_{2}$, and $\tilde\beta_{1}$ is uncorrelated with $\hat\beta_{2}$, so
$\Var(\hat\beta_{1})=\Var(\tilde\beta_{1})+a\Var(\hat\beta_{2})a\xT$. Then
\[
  MSE(\hat\beta_{1})-MSE(\tilde\beta_{1})=a\big[\underbrace{\Var(\hat\beta_{2})-\beta_{2}\beta_{2}\xT}_{\text{gap}}\big]a\xT .
\]
\Ans{Omitting $X_{2}$ wins in MSE when the gap is p.s.d.: $\beta_{2}$ is small relative to its sampling error. With one omitted variable,
$|\beta_{2}|<SD(\hat\beta_{2})$ (true $t$-ratio below 1).}
\begin{why}\Why omitting $X_{2}$ saves the variance of estimating $\beta_{2}$ but pays the squared bias $\beta_{2}^{2}$; both reach $\beta_{1}$ through $a$, so only their difference matters.\end{why}
\step{(d)} $\hat\beta_{1}=\beta_{1}+(X_{1}\xT M_{2}X_{1})^{-1}X_{1}\xT M_{2}u$: unbiased, variance $\sigma^{2}(X_{1}\xT M_{2}X_{1})^{-1}\ge\sigma^{2}(X_{1}\xT X_{1})^{-1}$.
Also $e=M_{X}y=M_{X}u$ (as $M_{X}X_{1}=O$), so $\E(e\xT e)=\sigma^{2}(n-k)$.
\Ans{Irrelevant variables: no bias, lost efficiency. Omitted relevant variables: bias, smaller variance.}
\end{solbox}

\begin{qbox}{D7.\ GLS}{\likely\ \newtag}
$y=X\beta+u$, $\E(u)=0$, $\Var(u)=\sigma^{2}\Omega$, $\Omega$ known p.d.
\begin{flatlist}
  \item[(a)] Show OLS is unbiased; find its true variance. Why are the usual OLS standard errors wrong?
  \item[(b)] Let $\Omega^{-1}=T\xT T$. Transform by $T$, show the errors are spherical, derive $\hat\beta_{GLS}$ and its variance.
  \item[(c)] Show $\Var(b_{OLS})-\Var(\hat\beta_{GLS})$ is p.s.d.
  \item[(d)] When does OLS equal GLS?
\end{flatlist}
\end{qbox}
\begin{solbox}
\step{(a)} $b=\beta+(X\xT X)^{-1}X\xT u$, $\E(b)=\beta$ (only $\E(u)=0$ needed). $\Var(b)=\sigma^{2}(X\xT X)^{-1}X\xT\Omega X(X\xT X)^{-1}\ne\sigma^{2}(X\xT X)^{-1}$; $s^{2}$ is biased too,
so reported standard errors and tests are wrong.
\step{(b)} $\Omega^{-1}=T\xT T\Rightarrow T\Omega T\xT=I$. With $y_{*}=Ty$, $X_{*}=TX$, $u_{*}=Tu$: $\Var(u_{*})=\sigma^{2}T\Omega T\xT=\sigma^{2}I$. OLS on the transformed model:
$\hat\beta_{GLS}=(X\xT\Omega^{-1}X)^{-1}X\xT\Omega^{-1}y$, $\Var(\hat\beta_{GLS})=\sigma^{2}(X\xT\Omega^{-1}X)^{-1}$.
\step{(c)} The transformed model satisfies the classical assumptions, and $b_{OLS}=(X\xT X)^{-1}X\xT T^{-1}y_{*}$ is linear and unbiased in $y_{*}$; Gauss--Markov (D3) gives p.s.d.
\step{(d)} $\Omega=I$, or more generally $\Omega X=XK$ with $K$ nonsingular: then $\Omega^{-1}X=XK^{-1}$ and
$\hat\beta_{GLS}=(K^{-1\prime}X\xT X)^{-1}K^{-1\prime}X\xT y=(X\xT X)^{-1}X\xT y$.
\begin{why}\Why GLS reweights so noisy or correlated observations count less; it is OLS on data rescaled to have spherical errors.\end{why}
\end{solbox}

\begin{qbox}{D8.\ Heteroskedasticity and WLS}{\opt\ \newtag}
$y_{i}=x_{i}\xT\beta+u_{i}$, $\Var(u_{i})=\sigma^{2}z_{i}^{2}$, $z_{i}$ observed.
\begin{flatlist}
  \item[(a)] Write $\Omega$ and $T$. Show WLS $=$ OLS of $y_{i}/z_{i}$ on $x_{i}/z_{i}$.
  \item[(b)] If $x_{i}=(1,z_{i})\xT$, write the transformed regression. What happens to the intercept?
\end{flatlist}
\end{qbox}
\begin{solbox}
\step{(a)} $\Omega=\mathrm{diag}(z_{i}^{2})$, $T=\mathrm{diag}(1/z_{i})$, $T\xT T=\Omega^{-1}$. The transformed model $\frac{y_{i}}{z_{i}}=\frac{x_{i}\xT}{z_{i}}\beta+\frac{u_{i}}{z_{i}}$ has
$\Var(u_{i}/z_{i})=\sigma^{2}$: OLS on it is GLS.
\step{(b)} $y_{i}=\beta_{1}+\beta_{2}z_{i}+u_{i}$ becomes $\frac{y_{i}}{z_{i}}=\beta_{1}\frac{1}{z_{i}}+\beta_{2}+\frac{u_{i}}{z_{i}}$.
\Ans{The original slope $\beta_{2}$ becomes the intercept; the original intercept $\beta_{1}$ becomes the coefficient on $1/z_{i}$. Do not add another constant.}
\end{solbox}

\begin{qbox}{D9.\ AR(1) errors}{\must\ \newtag}
$y_{t}=x_{t}\xT\beta+u_{t}$, $u_{t}=\rho u_{t-1}+\varepsilon_{t}$, $\varepsilon_{t}\sim IID(0,\sigma_{\varepsilon}^{2})$, $|\rho|<1$.
\begin{flatlist}
  \item[(a)] Write $u_{t}$ as an MA($\infty$). Derive $\Var(u_{t})$, $\Cov(u_{t},u_{t-s})$ and $\Omega$.
  \item[(b)] With $\rho$ known: give the quasi-differenced model for $t=2,\dots,n$ and the treatment of the first observation.
  \item[(c)] Write $T$ (Prais--Winsten) for $n=3$ and verify $T\xT T=(1-\rho^{2})\Omega^{-1}$. What does dropping the first row lose?
  \item[(d)] Repeat (b) for AR(2) with $\rho_{1},\rho_{2}$ known.
  \item[(e)] Now $\rho$ is \emph{unknown}. Write the model in two ways (linear in $\beta$ given $\rho$; linear in $\rho$ given $\beta$) and describe
        the Cochrane--Orcutt iteration, including the starting value and stopping rule. Why is OLS a legitimate start?
\end{flatlist}
\end{qbox}
\begin{solbox}
\step{(a)} $(1-\rho L)u_{t}=\varepsilon_{t}\Rightarrow u_{t}=\sum_{j\ge0}\rho^{j}\varepsilon_{t-j}$ (so $\E(u_{t})=0$). The $\varepsilon$'s are uncorrelated:
$\Var(u_{t})=\frac{\sigma_{\varepsilon}^{2}}{1-\rho^{2}}\equiv\sigma_{u}^{2}$, $\Cov(u_{t},u_{t-s})=\rho^{s}\sigma_{u}^{2}$, so
$\Var(u)=\sigma_{u}^{2}\Omega$ with $\Omega_{ts}=\rho^{|t-s|}$.
\begin{why}\Why $u_{t}$ and $u_{t-s}$ share the shocks $\varepsilon_{t-s},\varepsilon_{t-s-1},\dots$; in $u_{t}$ they carry the extra weight $\rho^{s}$.\end{why}
\step{(b)} For $t\ge2$: $y_{t}-\rho y_{t-1}=(x_{t}-\rho x_{t-1})\xT\beta+\varepsilon_{t}$. For $t=1$: multiply by $\sqrt{1-\rho^{2}}$; the error $\sqrt{1-\rho^{2}}\,u_{1}$ has variance
$(1-\rho^{2})\sigma_{u}^{2}=\sigma_{\varepsilon}^{2}$, equal to the others. OLS on all $n$ transformed rows is exact GLS.
\step{(c)}
\[
  T=\begin{bmatrix}\sqrt{1-\rho^{2}}&0&0\\-\rho&1&0\\0&-\rho&1\end{bmatrix},\quad
  T\xT T=\begin{bmatrix}1&-\rho&0\\-\rho&1+\rho^{2}&-\rho\\0&-\rho&1\end{bmatrix}=(1-\rho^{2})\Omega^{-1}.
\]
The check is $\Var(Tu)=\sigma_{\varepsilon}^{2}I$: equal diagonal, zero off-diagonal. Dropping row 1 (Cochrane--Orcutt) makes the $(1,1)$ entry $\rho^{2}$ instead of $1$: no longer
proportional to $\Omega^{-1}$, so not exact GLS (a small efficiency loss).
\step{(d)} $y_{t}-\rho_{1}y_{t-1}-\rho_{2}y_{t-2}=(x_{t}-\rho_{1}x_{t-1}-\rho_{2}x_{t-2})\xT\beta+\varepsilon_{t}$, $t=3,\dots,n$; GLS apart from the first two observations (step 1 of A4).
\step{(e) Two forms of the same model.}
\[
  \text{(A)}\ \ y_{t}-\rho y_{t-1}=(x_{t}-\rho x_{t-1})\xT\beta+\varepsilon_{t}\quad(\text{linear in }\beta\text{ given }\rho),
\]
\[
  \text{(B)}\ \ u_{t}=\rho u_{t-1}+\varepsilon_{t},\ u_{t}=y_{t}-x_{t}\xT\beta\quad(\text{linear in }\rho\text{ given }\beta).
\]
\begin{flatnum}
  \item Start: OLS of $y$ on $x$ gives $\hat\beta^{(0)}$ and $\hat u_{t}=y_{t}-x_{t}\xT\hat\beta^{(0)}$.
  \item (B): OLS of $\hat u_{t}$ on $\hat u_{t-1}$ (no constant): $\hat\rho^{(1)}=\sum\hat u_{t}\hat u_{t-1}/\sum\hat u_{t-1}^{2}$.
  \item (A): quasi-difference with $\hat\rho^{(1)}$, OLS gives $\hat\beta^{(1)}$; new residuals from the \emph{original} data $\hat u_{t}=y_{t}-x_{t}\xT\hat\beta^{(1)}$.
  \item Repeat 2--3; stop when $|\hat\rho^{(j)}-\hat\rho^{(j-1)}|$, $|\hat\beta^{(j)}-\hat\beta^{(j-1)}|$ and $|SSR^{(j)}-SSR^{(j-1)}|$ are all below a tolerance.
\end{flatnum}
\begin{why}\Why (A) needs $\rho$ to build the data, (B) needs $\beta$ to build the errors; alternating solves both. OLS is a legitimate start because it needs only $\E(u)=0$ to be unbiased,
so its residuals are a sensible first estimate of $u_{t}$. The result is feasible GLS. For AR(2), (B) regresses $\hat u_{t}$ on $\hat u_{t-1},\hat u_{t-2}$ (A4).\end{why}
\end{solbox}

\begin{qbox}{D10.\ Instrumental variables and 2SLS}{\likely\ \newtag}
$y=X\beta+u$, $X$ is $n\times k$, $\E(u)=0$, $\Var(u)=\sigma^{2}I$, but $\E(X\xT u)\ne0$. Instruments $W$ ($n\times l$, $l\ge k$) satisfy $\E(W\xT u)=0$.
\begin{flatlist}
  \item[(a)] Write $b-\beta$ for OLS and explain why OLS is biased here.
  \item[(b)] For $l=k$, derive $b_{IV}$ from the sample moment condition $W\xT(y-X\beta)=0$.
  \item[(c)] For $l>k$, derive 2SLS (i) in two OLS stages and (ii) by minimizing $(y-X\beta)\xT P_{W}(y-X\beta)$. Show they coincide.
  \item[(d)] Show that 2SLS reduces to $b_{IV}$ when $l=k$.
\end{flatlist}
\end{qbox}
\begin{solbox}
\step{(a)} $b-\beta=(X\xT X)^{-1}X\xT u$ and $\E(X\xT u)\ne0$, so the error term has nonzero mean: biased.
\begin{why}\Why OLS forces residuals orthogonal to $X$; if the true errors are not, that forcing pushes $b$ away from $\beta$.\end{why}
\step{(b)} $W\xT(y-Xb_{IV})=0\Rightarrow W\xT Xb_{IV}=W\xT y\Rightarrow b_{IV}=(W\xT X)^{-1}W\xT y=\beta+(W\xT X)^{-1}W\xT u$ ($W\xT X$ square, invertible). The error now involves $W\xT u$, with mean zero.
\step{(c)(i)} Stage 1: $\hat X=P_{W}X$. Stage 2: $b_{2SLS}=(\hat X\xT\hat X)^{-1}\hat X\xT y=(X\xT P_{W}X)^{-1}X\xT P_{W}y$ (as $P_{W}P_{W}=P_{W}$).
\step{(c)(ii)} $Q(\beta)=y\xT P_{W}y-2\beta\xT X\xT P_{W}y+\beta\xT X\xT P_{W}X\beta$; FOC $-2X\xT P_{W}y+2X\xT P_{W}X\beta=0$ gives the same $\beta$.
\begin{why}\Why with $l>k$ there are more moment conditions than parameters, so $W\xT(y-X\beta)=0$ cannot hold exactly; 2SLS shrinks its instrument-space part $P_{W}(y-X\beta)$.
Equivalently it keeps only the part of $X$ the instruments explain, which is uncorrelated with $u$.\end{why}
\step{(d)} For $l=k$, $(X\xT P_{W}X)^{-1}=(W\xT X)^{-1}(W\xT W)(X\xT W)^{-1}$, so
$b_{2SLS}=(W\xT X)^{-1}(W\xT W)(X\xT W)^{-1}X\xT W(W\xT W)^{-1}W\xT y=(W\xT X)^{-1}W\xT y=b_{IV}$.
\end{solbox}

\topic{Removed as beyond Lessons 1--8 (revisit before the final)}
{\small Box--Cox model by maximum likelihood; GNR and LM tests for AR errors; ETM 3.16, 3.17, 4.23, 7.2, 7.5--7.7, 8.1.}

\begin{transcriptnotes}
  \item A1, A2: \texttt{614midf24.pdf} (Fall 2024; Q1 weight 30, Q3 weight 70; no Q2). A3, A4: Fall 2023 midterm (Q3 weight 40, Q1 weight 30); its Q2 is identical to A1.
  \item B: \texttt{hw01.tex}. C: ETM exercises 2.15, 2.19, 3.14, 3.15, 4.10, 4.12. D: new questions on lessons 01--08.
  \item Critical values from \texttt{scipy.stats}. A1 corrects two slips in the handwritten solution collection ($TSS=126.08$, and comparing $2.93$ with $3$).
\end{transcriptnotes}

\end{document}
